\SourceLeaf{319}%
bis auf die fünfte Ordnung (ausschl.) richtig sind, indem \(s\) wie eine Grösse erster Ordnung betrachtet wird, und dass man, ohne den Grad der Schärfe zu vermindern, in den eingeklammerten Gliedern rechter Hand statt \(q,\alpha l\) und \(u\) auch \(s\mathord{.}\cos U\), \(s\mathord{.}\frac{\sin U}{\cos Q}\), \(s\mathord{.}\sin U\mathord{.}\operatorname{tang}Q\) substituiren darf.

\Article{22}
Es müssen nun zuvörderst die Grössen \(B,b,T,t,r\), welche auf der Ellipsoidfläche ihre Bedeutung haben, mit ihren Correlaten auf der Kugelfläche \(Q,q,U,u,As\) verglichen werden. Alle dafür hier aufzustellenden Gleichungen werden bis wenigstens auf die dritte Ordnung (einschl.) genau sein, und, dass dieser Bedingung genügt werde, wird sich aus der Entwicklung selbst leicht erkennen lassen.

Wendet man die im 8. Art. gegebene Reihe auf unsere beiden Punkte an, so müssen die dort allgemein mit \(p\) und \(q\) bezeichneten Grössen nach unserer jetzigen Bezeichnung ausgedrückt werden
\begin{center}
\begin{tabular}{ll}
für den ersten Punkt durch & \(B+\frac12b-P\) und \(\frac12q\)\\
für den zweiten Punkt durch & \(B-\frac12b-P\) und \(-\frac12q\)
\end{tabular}
\end{center}
und wir haben demnach die beiden Gleichungen
\begingroup\fontsize{10}{12}\selectfont
\[
\begin{aligned}
B+\tfrac12b={}&P+\frac{\cos\theta}{2\cos\varphi}\mathord{.}q-\frac{3ee}{8\cos\varphi^2}\mathord{.}\cos P\mathord{.}\sin P\mathord{.}qq\\
&+\frac{ee}{16\cos\varphi^3\cos\theta}(-\cos P^2+\sin P^2+ee(5\cos P^2\mathord{.}\sin P^2-\sin P^4))q^3\\
B-\tfrac12b={}&P-\frac{\cos\theta}{2\cos\varphi}\mathord{.}q-\frac{3ee}{8\cos\varphi^2}\mathord{.}\cos P\mathord{.}\sin P\mathord{.}qq\\
&-\frac{ee}{16\cos\varphi^3\cos\theta}(-\cos P^2+\sin P^2+ee(5\cos P^2\mathord{.}\sin P^2-\sin P^4))q^3
\end{aligned}
\]
\endgroup
Durch Addition und Subtraction ergibt sich also
\[
\begin{aligned}
B={}&P-\frac{3ee}{8\cos\varphi^2}\mathord{.}\cos P\mathord{.}\sin P\mathord{.}qq\quad .\ .\ .\ .\ .\ .\ .\ .\ .\ .\ .\ .\ .\ .\ .\ .\ .\ .\ .\quad(4)\\
b={}&\frac{\cos\theta}{\cos\varphi}\mathord{.}q+\frac{ee}{8\cos\varphi^3\cos\theta}(-\cos P^2+\sin P^2\\
&\qquad+ee(5\cos P^2\mathord{.}\sin P^2-\sin P^4))q^3\qquad(5)
\end{aligned}
\]
Man sieht übrigens leicht, dass die Gleichung (4) um Grössen vierter, die Gleichung (5) hingegen nur um Grössen fünfter Ordnung ungenau ist.

Um \(T\) und \(t\) mit \(U\) und \(u\) zu vergleichen, werden die am Schluss des 15. Art. entwickelten Formeln benutzt werden müssen, denen eine Voraussetzung
